El informe documenta el desarrollo
de un sistema
recomendador de inversiones para FinPUC, una plataforma ficticia de asesor\'ia
financiera que genera portafolios personalizados para clientes del mercado
accionario estadounidense. El sistema atiende cinco perfiles de riesgo, desde
muy conservador hasta muy arriesgado, y opera bajo un doble objetivo:
maximizar el retorno esperado del cliente y maximizar la utilidad de la empresa
por comisiones de transacci\'on.

El universo de trabajo comprende 601 acciones, obtenidas tras filtrar 1\,406
tickers. Se implementaron y compararon tres capas
metodol\'ogicas de complejidad creciente. La primera es la optimizaci\'on
media-varianza de Markowitz, resuelta mediante programaci\'on cuadr\'atica con
CVXPY, con rebalanceo semestral cada 126 d\'ias h\'abiles. La segunda extiende
Markowitz con el modelo Black-Litterman, que combina un prior de equilibrio de
mercado con cuatro variantes de \textit{views} cuantitativas: momentum general,
momentum por capitalizaci\'on a 6 meses y a 1 a\~no, y un escenario macro de
desempleo asumido. La tercera capa es una simulaci\'on Monte Carlo de 5\,000
trayectorias durante 260 semanas que modela la din\'amica completa del cliente: aceptaci\'on de recomendaciones, retiro ante p\'erdidas y comisiones; para
las seis t\'ecnicas evaluadas. Como caso base se utiliza un benchmark
equiponderado de los 20 activos con mayor capitalizaci\'on, cuyo desempe\~no no
depende de estimaci\'on de par\'ametros.

Los resultados muestran que Black-Litterman con momentum por
capitalizaci\'on a 6 meses domina el ranking del recomendador, alcanzando una
riqueza terminal esperada de USD~8\,974 y un score P4 de 9\,417 en el perfil
muy arriesgado con datos de pandemia. El benchmark equiponderado se confirma
como un competidor exigente en retorno bruto (Sharpe~2.05), pero con
volatilidad y drawdown significativamente mayores. Incluir el per\'iodo de
pandemia (2020 - 2022) en el entrenamiento deteriora los portafolios agresivos
de Markowitz hasta en 1.1 puntos de Sharpe, mientras que las variantes
Black-Litterman muestran mayor estabilidad frente al cambio de r\'egimen.

\textbf{Palabras clave:} optimizaci\'on de portafolios, Markowitz,
Black-Litterman, simulaci\'on Monte Carlo, sistema recomendador, rebalanceo
semestral, CVXPY.